\SourcePage{129}{132}
\section{Addition of angular momenta}

Consider a system consisting of two weakly interacting parts. When the
interaction is entirely neglected, the conservation of angular momentum holds
for each part independently, so that the angular momenta $\mathbf L_1$ and
$\mathbf L_2$ of the two parts can be summed to give the total angular
momentum $\mathbf L$ of the system. In the next approximation, once the weak interaction is included,
the conservation laws for $\mathbf L_1$ and $\mathbf L_2$ cease to hold
rigorously; nonetheless, the quantum numbers $L_1$ and $L_2$, which fix the
squares of those angular momenta, remain “good” and can still serve for an
approximate characterization of the system's state. In a pictorial, classical treatment of the angular momenta, one may
say that in this approximation $\mathbf L_1$ and $\mathbf L_2$ precess about
the direction of $\mathbf L$ while their magnitudes remain unchanged.

The consideration of such systems raises the question of the
\emph{rule for adding angular momenta}. What values of $L$ are possible for
given $L_1$ and $L_2$? The addition rule for
\SourcePage{130}{133}
the angular-momentum projections is evident: from
$\hat L_z=\hat L_{1z}+\hat L_{2z}$ it follows that
\begin{equation}
  M=M_1+M_2.
  \tag{31.1}
\end{equation}
There is no equally simple relation for the operators of the squares of the
angular momenta, and we derive their “addition rule” as follows.

If the quantities $\mathbf L_1^2$, $\mathbf L_2^2$, $L_{1z}$, and $L_{2z}$
are chosen as a complete set of physical quantities\footnote{Together with a
number of other quantities which, along with the four stated here, make up a
complete set. These remaining quantities play no part in the following
argument; for conciseness we do not mention them at all and provisionally
call the system of the four stated quantities complete.}, each state is
specified by the numbers $L_1$, $L_2$, $M_1$, and $M_2$. For given $L_1$ and
$L_2$, the numbers $M_1$ and $M_2$ take, respectively, $2L_1+1$ and
$2L_2+1$ values, so there are altogether
$(2L_1+1)(2L_2+1)$ different states with the same $L_1$ and $L_2$. We denote
the wave functions of the states in this description by
$\varphi_{L_1L_2M_1M_2}$.

Instead of the four quantities just stated, one can choose the four quantities
$\mathbf L_1^2$, $\mathbf L_2^2$, $\mathbf L^2$, and $L_z$ as a complete
set. Each state is then characterized by the numbers $L_1$, $L_2$, $L$, and
$M$ (we denote the corresponding wave functions by $\psi_{L_1L_2LM}$).
For given $L_1$ and $L_2$ there must, of course, still be
$(2L_1+1)(2L_2+1)$ different states; that is, for fixed $L_1$ and $L_2$, the
pair $L,M$ can take $(2L_1+1)(2L_2+1)$ pairs of values. These values can be
found by the following argument.

Adding the different allowed values of $M_1$ and $M_2$ gives the corresponding
values of $M$:
\begin{flushleft}
\(
\begin{array}{c@{\qquad}c@{\qquad}c}
 M_1 & M_2 & M\\[2pt]
 \hline
 L_1 & L_2 & L_1+L_2\\[2pt]
 \multicolumn{2}{c}{
   \left.\begin{array}{cc}L_1&L_2-1\\L_1-1&L_2\end{array}\right\}}
 &L_1+L_2-1\\[8pt]
 \multicolumn{2}{c}{
   \left.\begin{array}{cc}L_1-1&L_2-1\\L_1&L_2-2\\L_1-2&L_2\end{array}\right\}}
 &L_1+L_2-2\\[-1pt]
 \multicolumn{3}{c}{\dotfill}
\end{array}
\)
\end{flushleft}
The largest possible value of $M$ is evidently $M=L_1+L_2$, and it
corresponds to one state $\varphi$ (one pair of values $M_1,M_2$).
Consequently, the largest possible $M$ in the states $\psi$, and hence the
largest $L$, is also $L_1+L_2$. Next, there are two states $\varphi$ with
$M=L_1+L_2-1$. There must therefore also be two states $\psi$ with this value
of $M$; one
\SourcePage{131}{134}
is a state with $L=L_1+L_2$ (and $M=L-1$), and the other has
$L=L_1+L_2-1$ (with $M=L$). For $M=L_1+L_2-2$ there are three different
states $\varphi$. This means that, in addition to $L=L_1+L_2$ and
$L=L_1+L_2-1$, the value $L=L_1+L_2-2$ is also possible.

This argument can be continued in the same way as long as decreasing $M$ by
one increases by one the number of states with that value of $M$. It is easy
to see that this continues until $M$ reaches $|L_1-L_2|$. As $M$ is reduced
further, the number of states ceases to increase and remains equal to
$2L_2+1$ (if $L_2\leq L_1$). Thus $|L_1-L_2|$ is the smallest possible
value of $L$.

We therefore obtain the result that, for given $L_1$ and $L_2$, $L$ can take
the values
\begin{equation}
  L=L_1+L_2,L_1+L_2-1,\ldots,|L_1-L_2|,
  \tag{31.2}
\end{equation}
altogether $2L_2+1$ different values when $L_2\leq L_1$. One readily checks
that these indeed give $(2L_1+1)(2L_2+1)$ different pairs of values $M,L$.
It is important to note that, apart from the $2L+1$ different values of $M$
for a given $L$, each possible value of $L$ in (31.2) corresponds to just one
state.

This result can be represented pictorially by the so-called
\emph{vector model}. Introduce two vectors $\mathbf L_1$ and $\mathbf L_2$
of lengths $L_1$ and $L_2$. The values of $L$ are represented by the
integer-valued lengths of vectors $\mathbf L$ obtained by vector addition of
$\mathbf L_1$ and $\mathbf L_2$; the largest value, $L_1+L_2$, is obtained
when $\mathbf L_1$ and $\mathbf L_2$ are parallel, and the smallest,
$|L_1-L_2|$, when they are antiparallel.

In states with definite angular momenta $L_1$, $L_2$, and definite total
angular momentum $L$, the scalar products $\mathbf L_1\mathbf L_2$,
$\mathbf L\mathbf L_1$, and $\mathbf L\mathbf L_2$ also have definite
values. These values are readily found. To calculate
$\mathbf L_1\mathbf L_2$, write
$\hat{\mathbf L}=\hat{\mathbf L}_1+\hat{\mathbf L}_2$, or, on squaring and
transposing terms,
\[
  2\hat{\mathbf L}_1\hat{\mathbf L}_2
  =\hat{\mathbf L}^{2}-\hat{\mathbf L}_1^{2}-\hat{\mathbf L}_2^{2}.
\]
Replacing the operators on the right by their eigenvalues gives the
eigenvalue of the operator on the left:
\begin{equation}
 \mathbf L_1\mathbf L_2=\frac12
 \{L(L+1)-L_1(L_1+1)-L_2(L_2+1)\}.
 \tag{31.3}
\end{equation}
Similarly, we find
\begin{equation}
 \mathbf L\mathbf L_1=\frac12
 \{L(L+1)+L_1(L_1+1)-L_2(L_2+1)\}.
 \tag{31.4}
\end{equation}
\SourcePage{132}{135}
Let us now establish the \emph{rule for combining parities}. For a system
that consists of two independent parts, the wave function $\Psi$ equals the
product of the individual wave functions $\Psi_1$ and $\Psi_2$. Hence, if the latter two
have the same parity (both change sign, or both remain unchanged, when all
coordinates change sign), the wave function of the whole system is even. If,
on the other hand, $\Psi_1$ and $\Psi_2$ have opposite parities, $\Psi$ is
odd. These statements can be written as
\begin{equation}
  P=P_1P_2,
  \tag{31.5}
\end{equation}
where $P$ is the parity of the whole system, and $P_1$ and $P_2$ are the
parities of its parts. This rule plainly generalizes directly to a system
consisting of any number of noninteracting parts.

In particular, for a system of particles in a centrally symmetric field
(with the interaction between the particles treated as weak), the parity of
the state of the whole system is
\begin{equation}
  P=(-1)^{l_1+l_2+\ldots}
  \tag{31.6}
\end{equation}
(see (30.7)). Note that the exponent involves the algebraic sum of the
particle angular momenta, a quantity that, generally speaking, does not
coincide with their “vector sum,” i.e., with the angular momentum $L$ of
the system.

If a closed system decays into parts under the action of forces within the
system itself, its total angular momentum and parity must be conserved. This
may make the decay impossible even when it is energetically possible.

For example, consider an atom in an even state with angular momentum $L=0$,
whose energy would permit it to break up into a free electron and an ion
occupying an odd state, also with $L=0$. It is easy to show that this
process is in fact impossible (one calls it \emph{forbidden}). Indeed, by
conservation of angular momentum the free electron would also have to possess
zero angular momentum and would therefore be in an even state
($P=(-1)^0=+1$). But the state ion$+$free electron would then be odd, whereas
the initial state of the atom was even.
